\documentclass[10pt,a4paper]{amsart}
\usepackage[margin=19mm]{geometry}
\usepackage{amsmath,amssymb,mathtools}
\usepackage[scr=rsfs,cal=euler]{mathalfa}
\usepackage{xfrac}
\usepackage[unicode,hidelinks,bookmarks=false]{hyperref}
\newcommand{\ie}{i.e.\ }
\newcommand{\cf}{cf.\ }
\newcommand{\resp}{respectively\ }
\newcommand{\Subsection}{\subsection}
\providecommand{\nobreakdash}{\nobreak\hbox{-}}
\setlength{\emergencystretch}{2em}
\multlinegap0pt
\usepackage{amssymb}
\usepackage{caption}
\usepackage{tikz}
\newtheorem{restatable}{Restatable}
\newtheorem{lemma}{Lemma}
\newtheorem{theorem}{Theorem}
\usepackage{cleveref}
\crefname{restatable}{Restatable}{Restatables}
\crefname{lemma}{Lemma}{Lemmas}
\crefname{theorem}{Theorem}{Theorems}
\newcommand{\homclass}[4]{\hom_{#1\to#2}(#3,#4)}
\newcommand{\ith}[1]{$#1^{\text{th}}$}
\title{Hoffman-London graphs: When paths minimize\\ $H$-colorings among trees}
\author{David Galvin, Phillip Marmorino, Emily McMillon, JD Nir, Amanda Redlich}
\date{Source: arXiv:2512.23828v1 (CC BY 4.0)}
\begin{document}
\maketitle

\noindent\emph{Source and licence.} Hoffman-London graphs: When paths minimize\\ $H$-colorings among trees. Version arXiv:2512.23828v1 (CC BY 4.0), distributed by arXiv under the Creative Commons Attribution 4.0 International licence, http://creativecommons.org/licenses/by/4.0/, as recorded in the arXiv licence metadata for this exact version and shown on the abstract page https://arxiv.org/abs/2512.23828v1. Full licence notice: \texttt{licenses/arxiv-2512.23828-CC-BY-4.0.txt}.\par\smallskip

\noindent\textit{Editorial scope.} Counted unit: the article's theorem thm:even\_path\_minimizer (source0.tex lines 678--680). The article's complete original proof of it is delivered with it (source0.tex lines 682--744, one \texttt{proof} environment of 63 lines). No mathematics is written, repaired or re-derived: the statement and the proof are the article's own lines, in the article's own notation, and no retained line is substituted or re-flowed.  Retained uncounted as the source-local context the proof needs: lines 130--133; lines 137--139; lines 275--279; lines 292--294.  Nothing was omitted from any retained span, so the package carries no omission record.  No retained line contains a citation, so the wrapper carries no bibliography and no bibliography entry.  No retained line contains a figure environment, an image-inclusion command or drawing code, so no image is delivered and the package declares no asset dependency.  Editorial formatting only, in addition: an AMS \texttt{amsart} preamble that restates the article's own declarations and macros used by the retained lines, namely the environments \texttt{restatable}, \texttt{lemma}, \texttt{theorem} on separate counters, together with the standard packages those lines need.  The retained environments print in this document as Restatable 1, Lemma 1, Theorem 1 in order of appearance, whereas the article prints the same items with the numbering of its own full document.  The complete native source of the article as distributed in the arXiv e-print for this exact version is bundled unchanged as \texttt{sources/191832/source0.tex}.
Our first contribution, captured by \cref{thm:ordering_framework}, is a test that can be applied to the automorphic similarity matrix of a graph $H$ which, when it applies, proves that $H$ is Hoffman-London. 
\begin{restatable}{theorem}{orderings} \label{thm:ordering_framework}
Let $H$ be a target graph, possibly with loops. Suppose there is an ordering of the automorphic similarity classes $H^1, \ldots, H^k$ of $H$ such that the automorphic similarity matrix $M$ has the increasing columns property. Then $H$ is Hoffman-London. 
\end{restatable}

\begin{restatable}{corollary}{uniquenessalt} \label{cor:uniquenessalt}
Let $H$ be a target graph that satisfies the hypotheses of \cref{thm:ordering_framework}. Suppose that for each $t \ge 2$, there exist two classes, $H^{a(t)}$ and $H^{b(t)}$, such that there is at least one $H$-coloring of $P_t$ that sends one endvertex of $P_t$ to a vertex in $H^{a(t)}$ and the other endvertex to a vertex in $H^{b(t)}$. Suppose further that for all $s \ge 2$, it holds that $\homclass{v}{b(t)}{P_s}{H} > \homclass{v}{a(t)}{P_s}{H}$, where $v$ is one of the end vertices of $P_s$. Then $H$ is strongly Hoffman-London. 
\end{restatable}

\begin{lemma} \label{lem:tree_walk}
Let $H$ be a target graph with automorphic similarity classes $H^1, \ldots, H^k$ and associated automorphic similarity matrix $M$. Let $T$ be a tree with root $v$.  Let $\mathbf{a}(H)$ be the row vector whose \ith{i} entry is $|H^i|$ and define $\mathbf{h}(T,v)$ to be the column vector whose \ith{i} entry is $\homclass{v}{i}{T}{H}$.

Then $\hom(T,H)={\mathbf a}(H) \mathbf{h}(T,v)$, and $\mathbf{h}(T,v)$ can be explicitly computed from $T$, $\mathbf{a}(H)$, and $M$, via recursion on $T$.
\end{lemma}

\begin{equation} \label{eq:treewalkalgorithm-pathcase} 
\mathbf{h}(T,v) = M \mathbf{h}(T-v,w)
\end{equation}

\begin{theorem} \label{thm:even_path_minimizer}
The path graph $P_{2m}$ is Hoffman-London for all $m \ge 1$ and strongly Hoffman-London for all $m \ge 2$.
\end{theorem}

\begin{proof}
Let $H = P_{2m}$ for some $m \ge 1$, specifically with vertices $u_1, \ldots, u_{2m}$ and edges $u_iu_{i+1}$ for $1 \le i \le 2m-1$. We claim that $P_{2m}$ admits exactly two automorphisms: the identity and the {\it reversal} map that sends $u_i$ to $u_{2m+1-i}$ for $i=1, \ldots, 2m$. To verify this claim, first note that an automorphism of $P_{2m}$ must preserve degrees and thus leaves must be mapped to leaves. If $u_1$ is mapped to $u_1$, then $u_2$, the unique neighbor of $u_1$, must be mapped to $u_2$, $u_3$ must be mapped to $u_3$, and so on, resulting in the identity map. If $u_1$ is instead mapped to $u_{2m}$, then $u_2$ must be mapped to $u_{2m-1}$, $u_3$ must be mapped to $u_{2m-2}$, and so on, giving resulting in the reversal map. It follows that $P_{2m}$ has $m$ automorphic similarity classes, each consisting of a pair of vertices $\{u_i,u_j\}$ such that $i+j=2m+1$. 

Set $H^i=\{u_i, u_{2m+1-i}\}$ for $i=1, \ldots, m$. The automorphic similarity matrix of $H$ with respect to this ordering of the automorphic similarity classes has the form 
\[ \begin{bmatrix}
0 & 1 & 0 & 0 & \cdots & 0 & 0 & 0\\
1 & 0 & 1 & 0 & \cdots & 0 & 0 & 0\\
0 & 1 & 0 & 1 & \cdots & 0 & 0 & 0\\
\vdots & \vdots & \vdots & \vdots & \ddots & \vdots & \vdots& \vdots\\
0 & 0 & 0 & 0 & \cdots & 1 & 0 & 1\\
0 & 0 & 0 & 0 & \cdots & 0 & 1 & 1
    \end{bmatrix}. \]
That is, it has $1$s on the first subdiagonal, $1$s on the first superdiagonal, a $1$ in the $(m,m)$ position, and $0$s everywhere else. It is straightforward to check that all terminal sums of the \ith{i} row are less than or equal to the corresponding terminal sums of the \ith{(i+1)} row for all $i < m$, and so this matrix satisfies the increasing columns property, establishing that $P_{2m}$ is Hoffman-London by \cref{thm:ordering_framework}.

We now turn to establishing that $P_{2m}$ is strongly Hoffman-London for $m \ge 2$. We seek to apply \cref{cor:uniquenessalt}. Towards this goal, we will show by induction on $t$ that
\begin{equation} \label{inq:paths}
\homclass{v}{1}{P_t}{P_{2m}} < \homclass{v}{2}{P_t}{P_{2m}} \le \homclass{v}{3}{P_t}{P_{2m}} \le \cdots \le \homclass{v}{m}{P_t}{P_{2m}},
\end{equation}
where $v$ is an endvertex of $P_t$.

When $t=2$, inequality \eqref{inq:paths} holds because $\homclass{v}{1}{P_2}{P_{2m}}=1$ and $\homclass{v}{i}{P_2}{P_{2m}}=2$ for $2 \le i \le m$. For $t > 2$, let $v'$ be the unique neighbor of $v$ in $P_t$. Using equation \eqref{eq:treewalkalgorithm-pathcase} from the proof of \cref{lem:tree_walk}, we have the recurrence relations
\begin{eqnarray*}
\homclass{v}{1}{P_t}{P_{2m}} & = & \homclass{v'}{2}{P_{t-1}}{P_{2m}}, \\
\homclass{v}{i}{P_t}{P_{2m}} & = & \homclass{v'}{i-1}{P_{t-1}}{P_{2m}} + \homclass{v'}{i+1}{P_{t-1}}{P_{2m}} \text{ when } 2 \le i \le m-1, \text{ and} \\
\homclass{v}{m}{P_t}{P_{2m}} & = & \homclass{v'}{m-1}{P_{t-1}}{P_{2m}} + \homclass{v'}{m}{P_{t-1}}{P_{2m}}.
\end{eqnarray*}
By induction, we have
\begin{equation}\label{inq:path_induction}
\homclass{v'}{1}{P_{t-1}}{P_{2m}} < \homclass{v'}{2}{P_{t-1}}{P_{2m}} \le \cdots \le \homclass{v'}{m}{P_{t-1}}{P_{2m}}.
\end{equation}
When $m=2$, we use the fact that $\homclass{v'}{1}{P_{t-1}}{P_{4}} > 0$ to say
\begin{eqnarray*}
\homclass{v}{1}{P_t}{P_{4}} &=& \homclass{v'}{2}{P_{t-1}}{P_{4}} \\
	&<& \homclass{v'}{1}{P_{t-1}}{P_{4}} + \homclass{v'}{2}{P_{t-1}}{P_{4}} \\
	&=& \homclass{v}{2}{P_t}{P_4}.
\end{eqnarray*}
which establishes inequality~\eqref{inq:paths}.

For $m \ge 3$, we first establish $\homclass{v}{1}{P_t}{P_{2m}} < \homclass{v}{2}{P_t}{P_{2m}}$. By inequality~\eqref{inq:path_induction}, we see that $\homclass{v'}{2}{P_{t-1}}{P_{2m}} \le \homclass{v'}{3}{P_{t-1}}{P_{2m}}$. Then, again using $\homclass{v'}{1}{P_{t-1}}{P_{2m}} > 0$, we have
\begin{eqnarray*}
\homclass{v}{1}{P_t}{P_{2m}} &=& \homclass{v'}{2}{P_{t-1}}{P_{2m}} \\
	&\le& \homclass{v'}{3}{P_{t-1}}{P_{2m}} \\
	&<& \homclass{v'}{1}{P_{t-1}}{P_{2m}} + \homclass{v'}{3}{P_{t-1}}{P_{2m}} \\
	&=& \homclass{v}{2}{P_t}{P_{2m}}.
\end{eqnarray*}
Next, we show $\homclass{v}{i}{P_t}{P_{2m}} \le \homclass{v}{i+1}{P_t}{P_{2m}}$ for $2 \le i \le m-2$ by using inequality~\eqref{inq:path_induction} to say
\begin{eqnarray*}
\homclass{v}{i}{P_t}{P_{2m}} &=& \homclass{v'}{i-1}{P_{t-1}}{P_{2m}} + \homclass{v'}{i+1}{P_{t-1}}{P_{2m}} \\
	&\le& \homclass{v'}{i}{P_{t-1}}{P_{2m}} + \homclass{v'}{i+2}{P_{t-1}}{P_{2m}} \\
	&=& \homclass{v}{i+1}{P_t}{P_{2m}}.
\end{eqnarray*}
Finally, we demonstrate $\homclass{v}{m-1}{P_t}{P_{2m}} \le \homclass{v}{m}{P_t}{P_{2m}}$ by once again using inequality~\eqref{inq:path_induction}:
\begin{eqnarray*}
\homclass{v}{m-1}{P_t}{P_{2m}} &=& \homclass{v'}{m-2}{P_{t-1}}{P_{2m}} + \homclass{v'}{m}{P_{t-1}}{P_{2m}} \\
	&\le& \homclass{v'}{m-1}{P_{t-1}}{P_{2m}} + \homclass{v'}{m}{P_{t-1}}{P_{2m}} \\
	&=& \homclass{v}{m}{P_t}{P_{2m}}.
\end{eqnarray*}
which establishes inequality~\eqref{inq:paths}.

Finally, we use inequality~\eqref{inq:paths} to show that $P_{2m}$ is strongly Hoffman-London. We consider the cases $m=2$ and $m > 2$ separately. For $m=2$ and $t \ge 2$, set $a(t)=1$ and $b(t)=2$. From inequality~\eqref{inq:paths}, we have $\homclass{v}{a(t)}{P_s}{P_{4}} < \homclass{v}{b(t)}{P_s}{P_{4}}$ for all $s \ge 2$. To apply \cref{cor:uniquenessalt} to show that $P_{4}$ is strongly Hoffman-London, it remains to show that for all $t \ge 2$ there is a $P_{4}$-coloring of $P_t$ that sends one endvertex $v$ of $P_t$ to some vertex in $H^1$ and the other endvertex $w$ to some vertex in $H^2$. For $t$ even, such a coloring is given by sending $v$ to $u_1\in H^1$, then sending the vertices of $P_t$ alternatively to $u_2$ and $u_1$, ending by sending $w$ to $u_2 \in H^2$. For $t$ odd, such a coloring is given by sending $v$ to $u_1\in H^1$, then sending the vertices of $P_t$ alternatively to $u_2$ and $u_1$ except for $w$, which is sent to $u_3 \in H^2$ instead of $u_1$.

For $m>2$, set $a(t)=1$ for $t \ge 2$, and set $b(t)=2$ for even $t \ge 2$ and set $b(t)=3$ for odd $t \ge 2$. From inequality~\eqref{inq:paths}, we again have $\homclass{v}{a(t)}{P_s}{P_{2m}} < \homclass{v}{b(t)}{P_s}{P_{2m}}$ for all $s \ge 2$. To exhibit a coloring of $P_t$ that sends $v$ to a vertex in $H^{a(t)}$ and $w$ to a vertex in $H^{b(t)}$, we proceed exactly as in the $m=2$ case: send $v$ to $u_1 \in H^1$ and proceed to alternate sending the vertices of $P_t$ to $u_2$ and $u_1$, ending by sending $w$ to $u_2\in H^2$ when $t$ is even and deviating by sending $w$ to $u_3 \in H^3$ instead of $u_1$ when $t$ is odd. 
\end{proof}

\end{document}
